
\documentclass{article}
\usepackage{colortbl}
\usepackage{makecell}
\usepackage{multirow}
\usepackage{supertabular}

\begin{document}

\newcounter{utterance}

\centering \large Interaction Transcript for game `eqbench', experiment `eq\_bench', episode 58 with reflexion{-}Qwen3.5{-}9B.
\vspace{24pt}

{ \footnotesize  \setcounter{utterance}{1}
\setlength{\tabcolsep}{0pt}
\begin{supertabular}{c@{$\;$}|p{.15\linewidth}@{}p{.15\linewidth}p{.15\linewidth}p{.15\linewidth}p{.15\linewidth}p{.15\linewidth}}
    \# & \multicolumn{2}{c}{Player} && \multicolumn{2}{c}{Game Master} \\
    \hline

    \theutterance \stepcounter{utterance}  
    & & & \multicolumn{4}{p{0.6\linewidth}}{
        \cellcolor[rgb]{0.9,0.9,0.9}{
            \makecell[{{p{\linewidth}}}]{
                \texttt{\tiny{[P1$\langle$GM]}}
                \texttt{You have attempted this same task before and did not succeed. These are the plans you wrote after those attempts.} \\
\texttt{Trial \#1: New plan: I failed because I assigned a high score to "Guilt" (8) despite the dialogue showing Tyler explicitly accusing Claudia of leaving him behind and treating her as an "afterthought." While Claudia denies being a burden, Tyler's aggressive questioning ("Is that why you're here now?") and the context of her past abandonment suggest her dominant emotion is defensive anger or hurt, not guilt. I should have analyzed the subtext of Tyler's accusations more deeply to identify "Resentment" as the primary driver of the conflict. Next time, I will re{-}evaluate the emotional weight of Tyler's final lines, recognizing that his confrontation stems from long{-}held bitterness. I will assign a high score (e.g., 9) to Resentment to reflect Tyler's feeling, and a moderate score (e.g., 4) to Guilt to acknowledge Claudia's defensive denial without overestimating her remorse. I will ensure the highest score aligns with the most prominent emotional theme revealed by the antagonist's dialogue.} \\
\\ 
\texttt{Trial \#2: New plan: I failed because I misattributed the emotions to the wrong character. The prompt asks for **Claudia's** emotional response, but I scored "Resentment" (9) and "Determination" (3), which are clearly Tyler's feelings based on his accusatory dialogue ("dirty work," "afterthought," "burden"). Claudia's lines ("I didn't mean it like that," "I try to understand," "You were my responsibility") indicate she is trying to explain herself and justify her past actions, which aligns with **Guilt**. I should have strictly identified the subject (Claudia) before scoring. Next time, I will analyze Claudia's specific dialogue to confirm her internal state: her defensive yet explanatory tone suggests she feels responsible for the pain she caused. I will assign a high score (e.g., 9) to **Guilt** to reflect her sense of responsibility and regret, and low scores (0{-}2) to the other options, ensuring the highest score matches the character's actual perspective.} \\
\\ 
\\ 
\\ 
\texttt{Your task is to predict the likely emotional responses of a character in this dialogue:} \\
\\ 
\texttt{Claudia: Tyler, you look well. Teaching suits you.} \\
\texttt{Tyler: Yeah, well, someone's gotta do the dirty work.} \\
\texttt{Claudia: I didn't mean it like that. I...} \\
\texttt{Tyler: I know what you meant. You always do, don't you?} \\
\texttt{Claudia: Not always, no. But I try to understand, at least.} \\
\texttt{Tyler: Understand? Like you understood when you left me alone?} \\
\texttt{Claudia: That's not fair, Ty. I had to...} \\
\texttt{Tyler: What? Chase your dreams? Leave your past behind?} \\
\texttt{Claudia: I never left you behind, Tyler. You were always...} \\
\texttt{Tyler: Always what? An afterthought? A burden?} \\
\texttt{Claudia: No, never a burden. You were my responsibility.} \\
\texttt{Tyler: Responsibility? Is that why you're here now? After all these years?} \\
\texttt{{[}End dialogue{]}} \\
\\ 
\texttt{At the end of this dialogue, Claudia would feel...} \\
\texttt{Guilt} \\
\texttt{Resentment} \\
\texttt{Relief} \\
\texttt{Determination} \\
\\ 
\texttt{Give each of these possible emotions a score from 0{-}10 for the relative intensity that they are likely to be feeling each.} \\
\\ 
\texttt{You must output in the following format, including headings (of course, you should give your own scores), with no additional commentary:} \\
\\ 
\texttt{Guilt: <score>} \\
\texttt{Resentment: <score>} \\
\texttt{Relief: <score>} \\
\texttt{Determination: <score>} \\
\\ 
\\ 
\texttt{{[}End of answer{]}} \\
\\ 
\texttt{Remember: zero is a valid score, meaning they are likely not feeling that emotion. You must score at least one emotion > 0.} \\
\\ 
\texttt{Your answer:} \\
            }
        }
    }
     \\ \\

    \theutterance \stepcounter{utterance}  
    & \multicolumn{4}{p{0.6\linewidth}}{
        \cellcolor[rgb]{0.561,0.737,0.561}{
            \makecell[{{p{\linewidth}}}]{
                \texttt{\tiny{[P1$\rangle$GM]}}
                \texttt{Guilt: 9} \\
\texttt{Resentment: 2} \\
\texttt{Relief: 1} \\
\texttt{Determination: 3} \\
            }
        }
    }
    & & \\ \\

    \theutterance \stepcounter{utterance}  
    & & & \multicolumn{2}{p{0.3\linewidth}}{
        \cellcolor[rgb]{0.9,0.9,0.9}{
            \makecell[{{p{\linewidth}}}]{
                \texttt{\tiny{[GM$|$GM]}}
                \texttt{Guilt: 9} \\
\texttt{Resentment: 2} \\
\texttt{Relief: 1} \\
\texttt{Determination: 3} \\
            }
        }
    }
    & & \\ \\

    \theutterance \stepcounter{utterance}  
    & & & \multicolumn{2}{p{0.3\linewidth}}{
        \cellcolor[rgb]{0.9,0.9,0.9}{
            \makecell[{{p{\linewidth}}}]{
                \texttt{\tiny{[GM$|$GM]}}
                \texttt{Guilt: 8} \\
\texttt{Resentment: 4} \\
\texttt{Relief: 0} \\
\texttt{Determination: 6} \\
            }
        }
    }
    & & \\ \\

    \theutterance \stepcounter{utterance}  
    & & & \multicolumn{2}{p{0.3\linewidth}}{
        \cellcolor[rgb]{0.9,0.9,0.9}{
            \makecell[{{p{\linewidth}}}]{
                \texttt{\tiny{[GM$|$GM]}}
                \texttt{game\_result = LOSE} \\
            }
        }
    }
    & & \\ \\

    \theutterance \stepcounter{utterance}  
    & & \multicolumn{4}{p{0.6\linewidth}}{
        \cellcolor[rgb]{0.95,0.95,0.95}{
            \makecell[{{p{\linewidth}}}]{
                \texttt{\tiny{[P1$\leftrightarrow$P]}}
                \texttt{Post{-}episode Reflexion (not sent to GM):} \\
\texttt{New plan: I failed because I assigned a high score to "Guilt" (9) despite the dialogue showing Tyler explicitly accusing Claudia of leaving him behind and treating her as an "afterthought." While Claudia denies being a burden, Tyler's aggressive questioning ("Is that why you're here now?") and the context of her past abandonment suggest her dominant emotion is defensive anger or hurt, not guilt. I should have analyzed the subtext of Tyler's accusations more deeply to identify "Resentment" as the primary driver of the conflict. Next time, I will re{-}evaluate the emotional weight of Tyler's final lines, recognizing that his confrontation stems from long{-}held bitterness. I will assign a high score (e.g., 9) to Resentment to reflect Tyler's feeling, and a moderate score (e.g., 4) to Guilt to acknowledge Claudia's defensive denial without overestimating her remorse. I will ensure the highest score aligns with the most prominent emotional theme revealed by the antagonist's dialogue.} \\
            }
        }
    }
    & \\ \\

\end{supertabular}
}

\end{document}
